\documentclass[../main.tex]{subfiles}

\begin{document}

\newcommand{\gpuModel}{RTX 3090}
\newcommand{\peakThroughputValue}{35.6}
\newcommand{\TFLOPS}{\ensuremath{\mathrm{TFLOP/s}}}
\newcommand{\GBS}{\ensuremath{\mathrm{GB/s}}}
\newcommand{\dramBandwidthValue}{935.8}
\newcommand{\peakThroughput}{\ensuremath{\peakThroughputValue\,\TFLOPS}}
\newcommand{\dramBandwidth}{\ensuremath{\dramBandwidthValue\,\GBS}}
\newcommand{\hardwareSpecifications}{%
  \begin{center}
    \begin{tabular}{l|rl}
      \gpuModel \\
      \hline
      Peak throughput & \peakThroughputValue & \TFLOPS\\
      DRAM bandwidth & \dramBandwidthValue & \GBS\\
    \end{tabular}
  \end{center}%
}

\subsection{Throughput, Arithmetic Intensity, and the Roofline Model}

\begin{frame}
  \frametitle{Driving Question}
  \hardwareSpecifications

  \begin{alertblock}{Is there a better performance metric?}
    We have been measuring our kernels with \textbf{latency}, but is there some
    other metric that better describes "performance"?
  \end{alertblock}

  \pause
  \begin{alertblock}{How fast \emph{can} we be?}
    Given our hardware constraints, how fast (theoretically) can we do our SGEMM operations?
  \end{alertblock}
\end{frame}

\begin{frame}
  \frametitle{Why latency alone isn't enough}
  \only<1>{
    \begin{itemize}
      \item Latency tells us how long one kernel launch takes, but it depends on
        the problem size.
    \end{itemize}
    \centering
    \includegraphics[width=0.7\textwidth]{generated/cpu_naive_tiled_latency.png}
  }

  \only<2>{
    \begin{itemize}
      \item To compare performance \textbf{across problem sizes},
        we need to account for both the work completed and the time required.
        This gives us \textbf{achieved throughput}.
    \end{itemize}
    \centering
    \includegraphics[width=0.66\textwidth]{generated/cpu_naive_tiled_throughput.png}
  }

  \only<3>{
    \begin{itemize}
      \item Since hardware is characterized by its \textbf{peak throughput}, the
        two values provide a first gauge of how much compute capability our
        kernel realizes.
    \end{itemize}
    \centering
    \includegraphics[width=0.7\textwidth]{generated/cpu_naive_tiled_peak_throughput.png}
  }

  \only<4>{
    \begin{itemize}
      \item Latency tells us how long one kernel launch takes, but it depends on
        the problem size.
      \item To compare performance \textbf{across problem sizes},
        we need to account for both the work completed and the time required.
        This gives us \textbf{achieved throughput}.
      \item Since hardware is characterized by its \textbf{peak throughput}, the
        two values provide a first gauge of how much compute capability our
        kernel realizes.
    \end{itemize}

    \begin{alertblock}{Why throughput}
      Throughput tells us
      \begin{enumerate}
        \item A general performance metric, no matter the problem size
        \item How well we are utilizing hardware
      \end{enumerate}
    \end{alertblock}
  }
\end{frame}

\begin{frame}
  \frametitle{Throughput as a performance measure}
  \begin{block}{Throughput}
    \textbf{Throughput} is the total useful work (FLOPs) done
    per unit time.
    \[
      \mathrm{Throughput} = \frac{\mathrm{FLOPs}}{\mathrm{time}}
    \]
  \end{block}

  \pause
  \begin{example}[Measuring Tiled Kernel]
    We observed $360$ms or $0.36$s of latency for \(n=8192\),
    where the total FLOPs required is \(2(8192)^3 \approx 1.10\) TFLOPs.
    Thus, we have an achieved throughput of
    \[\frac{1.10\mathrm{ TFLOPs}}{0.36\mathrm{s}} \approx 3.05\,\TFLOPS. \]

    Roughly $8.5$\% of peak throughput \peakThroughput.
  \end{example}
\end{frame}

\begin{frame}
  \frametitle{Characterizing our kernels data use}
  \begin{enumerate}
    \item Throughput characterizes how well our kernel utilizes the \textbf{compute hardware}.
    \item \pause One factor that impacts final throughput is how efficiently our kernels use data.
  \end{enumerate}

  \pause
  \begin{alertblock}{How can we quantify how "well" our kernel uses data?}
    We observed that we can improve the performance of our SGEMM kernel with
    tiling, an optimization that improves \emph{data use}.
    \pause \textbf{Can we quantify this data usage?}
  \end{alertblock}
\end{frame}

\begin{frame}
  \frametitle{Arithmetic Intensity: Characterizing our kernels}

  \begin{block}{Arithmetic Intensity}
    We define the \textbf{Arithmetic Intensity} (AI) of our kernel as the
    number of FLOPs done per bytes moved from DRAM.
    \[
      \mathrm{AI}=\frac{\mathrm{FLOPs}}{\mathrm{bytes\ moved}}
    \]
  \end{block}

  \pause
  \only<2>{
    \begin{example}[Naive kernel]

      For \(n,m,k=8192\), ignoring hardware L1 caching, we load (from DRAM)
      \[ n \times m \times 2k \times 4 \text{ bytes} \approx 4398 \text{ GB}.\]

      We do 1.10 TFLOPs of work, so we have an AI of
      \[ \mathrm{AI} = \frac{1.10 \text{ TFLOP}}{4.398 \text{ TB}} \approx 0.25\text{ FLOP/byte}.\]
    \end{example}
  }

  \only<3>{
    \begin{example}[Tiled kernel]
      For \(n, k, m = 8192\), in the worst case, our tiled kernel loads (from DRAM)
      \[ \frac{m}{32} \times \frac{n}{32} \times 2\left( \frac{k}{32}\right) \times 32 \times 32 \times 4 \text{ bytes} = 137.4\text{ GB}\]

      With a total work count of $1.10$ TFLOP, its arithmetic intensity is
      \[
        \mathrm{AI} = \frac{1.10 \text{ TFLOP}} {137.7\text{ GB}} \approx 8.0\ \text{FLOP/byte}.
      \]
    \end{example}
  }
\end{frame}

\begin{frame}
  \frametitle{The Memory Roof}
  \begin{block}{Memory Roof}
    Given the DRAM bandwidth and a kernel's AI, we say the \textbf{memory roof}
    of the kernel is
    \[\text{Memory Roof} = \text{AI} \times \text{DRAM bandwidth}.\]
  \end{block}

  \pause
  \begin{example}[Intuition]
    In order for the CUDA Cores to do work, we have
    to first move the data from DRAM to the SM's. \pause If the kernel only does a
    small amount of work on the data (low AI), it will likely finish before the
    next batch of data arrives. \pause \textbf{In this case, the final measured
    throughput will be constrained by the DRAM bandwidth.} We are \textbf{memory bound.}
  \end{example}

\end{frame}

\begin{frame}
  \frametitle{The Memory Roof}
  \hardwareSpecifications

  \begin{example}[Tiled Kernel]
    Our tiling kernel has AI of 8.0 FLOP/byte, meaning it does 8 floating point
    operations for every byte of data moved from DRAM.

    At \dramBandwidth{} of bandwidth, our \textbf{memory roof} is
    \[\dramBandwidth \times 8.0\text{ FLOP/byte} \approx 7.49\,\TFLOPS.\]

    \pause This is much lower than the peak throughput of \peakThroughput.
    \pause \textbf{No matter how we optimize this kernel, without increasing AI, we will never achieve greater than \(7.49\,\TFLOPS\)}.
  \end{example}
\end{frame}

\begin{frame}
  \frametitle{Hypothetical Kernel}
  \hardwareSpecifications
  \begin{example}
    Imagine we have some kernel that has double the AI at 16.0 FLOP/s.
    \pause Our memory roof is now
    \[\dramBandwidth \times 16.0\text{ FLOP/byte} \approx 15.0\,\TFLOPS.\]

    \pause Since \(15.0\,\TFLOPS < \peakThroughput\), this kernel will still be memory bound.
  \end{example}

\end{frame}

\begin{frame}
  \frametitle{Hypothetical Kernels}
  \hardwareSpecifications
  \begin{center}
    \begin{tabular}{c|c}
      AI (FLOP/byte) & Memory Roof (\TFLOPS) \\ \hline
      8  & 7.49 \\
      16 & 14.97 \\
      32 & 29.95 \\
      36 & 33.69 \\
      \alt<2->{\cellcolor{red!20}38}{38} &
      \alt<2->{\cellcolor{red!20}35.56}{35.56} \\
      \alt<2->{\cellcolor{red!20}40}{40} &
      \alt<2->{\cellcolor{red!20}37.43}{37.43}
    \end{tabular}
  \end{center}

  \pause
  Wait! Our peak throughput is \peakThroughput! At \(\text{AI} \ge 38\), we are
  \textbf{compute bound}. Our data can finally keep up with the CUDA Cores!
\end{frame}

\begin{frame}
  \frametitle{Ridge Point}
  \hardwareSpecifications

  \(\text{AI} = 38\text{ FLOP/byte}\) is known as the \textbf{ridge point} \(RP\),
  and can be found by solving:
  \begin{align*}
    \text{Memory Roof} &= \text{Peak throughput} \\
    \text{DRAM Bandwidth} \times \text{RP} &= \text{Peak throughput} \\
    \text{RP} &= \frac{\text{Peak throughput}}{\text{DRAM Bandwidth}} \\
    \text{RP} &= \frac{\peakThroughput}{\dramBandwidth} \approx 38.0\text{ FLOP/s}
  \end{align*}
\end{frame}

% \begin{frame}
%   \frametitle{Arithmetic Intensity (cont.)}

%   \begin{example}
%     Our tiled kernel has \(\text{AI} = 8.0 \text{ FLOP/byte}\).
%     Our DRAM bandwidth is \dramBandwidth.

%     \pause Our upper bound is
%     \(8.0 \text{ FLOP/byte} \times \dramBandwidth \approx 7.5\,\TFLOPS.\)

%     \pause Since this is less than peak throughput of \peakThroughput, we are memory bound and
%     would benefit from increasing AI.
%   \end{example}
% \end{frame}

\begin{frame}
  \frametitle{Roofline Model}
  \begin{block}{Roofline Model}
    For a given system, we can predict a the \textbf{upper bound} of throughput of a
    kernel given its AI.
    \pause Up to the \textbf{ridge point}, we are \textbf{memory bound},
    increasing AI will raise the upper bound of throughput.
    \pause Past that, we are \textbf{compute bound}, increasing AI
    will not raise the upper bound of throughput.
  \end{block}

  \onslide<1->
  \centering
  \only<1>{
    \includegraphics[height=0.54\textheight]{roofline_abstract/roofline_1.png}
  }
  \only<2>{
    \includegraphics[height=0.54\textheight]{roofline_abstract/roofline_2.png}
  }
  \only<3->{
    \includegraphics[height=0.54\textheight]{roofline_abstract/roofline_3.png}
  }

\end{frame}

\begin{frame}
  \frametitle{Roofline Model of our \gpuModel}
  \centering
  \includegraphics[width=0.9\textwidth]{generated/rtx3090_roofline.png}

  \centering
  Since we are memory bound, increasing AI should increase performance.
\end{frame}

\end{document}
